\section{第一节 函数}
\pt{求分段复合函数的解析式}
\begin{Problem}
    \difficulty{4}
    设
    $
    f(x)=
    \begin{cases}
    \sin{x}&,  |x|\leq \frac{\pi}{2} \\
    x&,  |x| > \frac{\pi}{2}
    \end{cases}
    $，
    $
    \phi(x)=
    \begin{cases}
    \arcsin x,&|x| \le 1 \\
    x,  &|x| > 1
    \end{cases}
    $
    则$f[\phi(x)]=\underline{\quad\quad\quad}$。
    \exercisesource{2025·660·第1题}
\end{Problem}
\begin{Answer}
    \multiplesolution{解一：代数法}
    \[
    f[\varphi(x)]=
    \begin{cases}
    \sin\varphi(x), & |\varphi(x)|\le \frac{\pi}{2} \\
    \varphi(x), & |\varphi(x)|> \frac{\pi}{2}
    \end{cases}
    \]
    \begin{enumerate}[label=(\arabic*)]
    \item 讨论 \(|\varphi(x)|\le \dfrac{\pi}{2}\)
        \begin{enumerate}[label=(\alph*)]
        \item 当 \(|x|\le 1\) 时，\(|\varphi(x)|=|\arcsin x|\le \dfrac{\pi}{2}\)，恒成立。
        \item 当 \(|x|> 1\) 时，\(|\varphi(x)|=|x|\le \frac{\pi}{2} \Rightarrow |x| > \frac{\pi}{2}\)。
        \end{enumerate}
        因此 \(|\varphi(x)| \le \frac{\pi}{2}\)时，$f[\phi(x)]=\sin[\phi(x)]=
        \begin{cases}
        x&,  |x|\le 1 \\
        \sin x&,  1<|x|\le \frac{\pi}{2}
        \end{cases}
        $。
    \item 讨论 \(|\varphi(x)|> \dfrac{\pi}{2}\)\begin{enumerate}[label=(\alph*)]
        \item 当 \(|x|\le 1\) 时，\(|\varphi(x)|=|\arcsin x|> \frac{\pi}{2}\)，无解。
        \item 当 \(|x|> 1\) 时，\(|\varphi(x)|=|x|> \dfrac{\pi}{2}\)，即 \(|x|> \frac{\pi}{2}\)。
        \end{enumerate}因此当$\varphi(x)> \frac{\pi}{2}$时，$f[\phi(x)]=\phi(x)=x,|x|>\frac{\pi}{2}$。
    \end{enumerate}   
    综上所述：
    \[
    f[\varphi(x)]=
    \begin{cases}
    x, & |x|\le 1 \text{ 或 } |x|> \dfrac{\pi}{2} \\
    \sin x, & 1<|x|\le \dfrac{\pi}{2}
    \end{cases}
    \]
    \multiplesolution{解二：图形法}
    $\varphi(x)$的图像，如图\ref{fig:ch01-1-2-tx-01-01}所示。因此有：
    \begin{figure}[htbp]  % [H] 表示图片固定在当前位置，也可用[htbp]自动排版
    \centering  % 核心：设置图片居中
    \includegraphics[width=0.8\textwidth]{figures/ch01/1-2/tx-01-01.pdf}
    \caption{分段函数 $\varphi(x)$ 图像}  % 设置标题
    \label{fig:ch01-1-2-tx-01-01}  % 设置标签，用于后续引用（标签名自定义，建议以fig:开头）
    \end{figure}
    \[
    \begin{aligned}
    f[\varphi(x)]
    &=
    \begin{cases}
    \sin\varphi(x), & |\varphi(x)| \le \dfrac{\pi}{2} \\
    \varphi(x),     & |\varphi(x)| > \dfrac{\pi}{2}
    \end{cases}
    \\
    &=
    \begin{cases}
    \sin\varphi(x), & |x| \le 1 \\
    \sin\varphi(x), & 1 < |x| \le \dfrac{\pi}{2} \\
    \varphi(x),     & |x| > \dfrac{\pi}{2}
    \end{cases}
    \\
    &=
    \begin{cases}
    x,             & |x| \le 1 \text{ 或 } |x| > \dfrac{\pi}{2} \\
    \sin x,        & 1 < |x| \le \dfrac{\pi}{2}
    \end{cases}
    \end{aligned}
    \]
\end{Answer}
\begin{Problem}
    \difficulty{3}
    设
    \(
    f(x)=\begin{cases}
    2x&, x\ge 0 \\
    x^2&,  x<0
    \end{cases},
    \quad
    g(x)=\begin{cases}
    x-1&,  x\ge 0 \\
    e^x&,  x<0
    \end{cases},
    \)
    求 \(f[g(x)]\)。
    \exercisesource{2025·660·第1题·举一反三}
    \exercisetip{内层函数不等式；带入内层函数时需注意}
\end{Problem}
\begin{Answer}
    \multiplesolution{解一：代数法}
    \[
    f[g(x)]=
    \begin{cases}
    2[g(x)]&, g(x)\ge 0 \\
    [g(x)]^2&,  g(x)<0
    \end{cases}
    \]
    \begin{enumerate}[label=(\arabic*)]
    \item 讨论 \(g(x)\ge 0\)的情况
        \begin{enumerate}[label=(\alph*)]
        \item 当 \(x \ge 0\) 时，\([g(x)]=(x-1)\ge 0 \Rightarrow x\ge 1\)。
        \item 当 \(x < 0\) 时，\([g(x)]=e^x\ge 0 \Rightarrow x <0\)。
        \end{enumerate}
        因此 \(g(x)\ge 0\)时，$f[g(x)]= 2[g(x)]=\begin{cases}
            2(x-1)&,x\ge 1\\
            2e^x&,x<0\\
            \end{cases}
        $。
    \item 讨论 \(g(x)< 0\)的情况
        \begin{enumerate}[label=(\alph*)]
        \item 当 \(x\ge 0\) 时，\([g(x)]=(x-1)< 0 \Rightarrow 0 \ge x <1\)。
        \item 当 \(x< 0\) 时，\([g(x)]=e^{x}<0\)，无解。
        \end{enumerate}
        因此\(g(x)< 0\)时，$f[g(x)]=[g(x)]^2=(x-1)^{2},0 \leq x<1$。
    \end{enumerate}   
    综上所述：
    \[
    f[g(x)]=
    \begin{cases}
    2(x-1)&, x\ge 1 \\
    (x-1)^{2}&,  0 \leq x<1\\
    2e^{x}&,  x<0
    \end{cases}
    \]
    \multiplesolution{解二：图形法}
    $g(x)$的图像，如图\ref{fig:ch01-1-2-tx-01-02-01}所示。因此有：
    \[
        \begin{aligned}
        f[g(x)]
        &=
        \begin{cases}
        2(g(x))&, g(x)\ge 0 \\
        (g(x))^{2}&,  g(x)<0
        \end{cases}
        \\
        &=
        \begin{cases}
        2(x-1)&, x\ge 1 \\
        (x-1)^{2}&,  0 \leq x<1\\
        2e^{x}&,  x<0
        \end{cases}
        \end{aligned}
    \]
    \begin{figure}[h]  % [H] 表示图片固定在当前位置，也可用[htbp]自动排版
    \centering  % 核心：设置图片居中
    \includegraphics[width=0.8\textwidth]{figures/ch01/1-2/tx-01-02-01.pdf}
    \caption{分段函数 $g(x)$ 图像}  % 设置标题
    \label{fig:ch01-1-2-tx-01-02-01}  % 设置标签，用于后续引用（标签名自定义，建议以fig:开头）
    \end{figure}
\end{Answer}
\begin{Problem}
    \difficulty{3}
    设
    \(
    f(x)=\begin{cases}
    (x-1)^2 & ,x\le 1 \\
    \dfrac{1}{1-x} & ,x>1
    \end{cases},
    \quad
    g(x)=\begin{cases}
    2x & ,x>0 \\
    x+2 & ,x\le 0
    \end{cases},
    \)
    求 $f[g(x)]$。
    \exercisesource{B站视频：BV1vx4y157hg}
\end{Problem}
\begin{Answer}
    \multiplesolution{解一：代数法}
    \[
    f[g(x)]=
    \begin{cases}
    [g(x)-1]^2&, g(x)\le 1 \\
    \dfrac{1}{1-g(x)}&, g(x)>1
    \end{cases}
    \]
    \begin{enumerate}[label=(\arabic*)]
    \item 讨论 \(g(x)\le 1\)的情况
        \begin{enumerate}[label=(\alph*)]
            \item 当 \(x > 0\) 时，\(g(x)=2x\le 1 \Rightarrow 0 < x \le \dfrac{1}{2}\)。
            \item 当 \(x \le 0\) 时，\(g(x)=2+x\le 1 \Rightarrow x\le -1\)。
        \end{enumerate}
        因此 \(g(x)\le 1\)时，$f[g(x)]=[g(x)-1]^2=\begin{cases}
            (2x-1)^2&,0 < x\le 1\\
            (x+1)^2&,x \le -1\\
        \end{cases} $。
    \item 讨论 \(g(x)> 1\)的情况
        \begin{enumerate}[label=(\alph*)]
            \item 当 \(x > 0\) 时，\(g(x)=2x> 1 \Rightarrow x>\frac{1}{2}\)。
            \item 当 \(x \le 0\) 时，\(g(x)=x+2> 1 \Rightarrow -1 < x \le 0\)。   
        \end{enumerate}
        因此\(g(x)> 1\)时，$f[g(x)]=\dfrac{1}{1-g(x)}=
        \begin{cases}
            \dfrac{1}{1-2x}&,x >\dfrac{1}{2}\\
            \dfrac{-1}{1+x}&,-1 < x \le 0\\
        \end{cases}$。
    \end{enumerate}   
    \multiplesolution{解二：图形法}
    $g(x)$的图像，如图\ref{fig:ch01-1-2-tx-01-03-01}所示。因此有：
    \begin{align*}
        &f[g(x)]=\begin{cases}
        [g(x)-1]^2&,\ g(x)\le 1 \\
        \dfrac{1}{1-g(x)}&,\ g(x)>1
        \end{cases}\\
        &=\begin{cases}
            (2x-1)^2&,\ 0 < x\le \dfrac{1}{2}\\
            (x+1)^2&,\ x \le -1\\
            \dfrac{1}{1-2x}&,\ x >\dfrac{1}{2}\\[0.5em]
            \dfrac{-1}{1+x}&,\ -1 < x \le 0\\
        \end{cases}
    \end{align*}
    \begin{figure}[htbp]
        \centering
        \includegraphics[width=0.8\textwidth]{ch01/1-2/tx-01-03-01.pdf}
        \caption{【3】题-分段函数 $g(x)$ 图像}
        \label{fig:ch01-1-2-tx-01-03-01}
    \end{figure}
\end{Answer}
\begin{Problem}
    \difficulty{2}
    已知$
    f(x)=\begin{cases}
        1+x&,\ x<0\\
        1&,\ x\ge 0
    \end{cases}
    $
    求$f[f(x)]$
    \exercisesource{B站视频：BV1mj411T7Tv}
\end{Problem}
\begin{Answer}
    \multiplesolution{解一：代数法}
    \[
    f[f(x)]=
    \begin{cases}
        1+f(x)&,\ f(x)<0\\
        1&,\ x\ge 0
    \end{cases}
    \]
    \begin{enumerate}[label=(\arabic*)]
        \item 考虑 \(f(x) < 0\) 的情况\begin{enumerate}[label=(\alph*)]
            \item 当 \(x < 0\) 时，\(f(x)=1+x< 0 \Rightarrow  x <-1\)。
            \item 当 $x \ge 0$时，$f(x)=1<0$，无解。
        \end{enumerate}$f[f(x)]=1+f(x)=2+x,\ x <-1$
        \item 考虑\(f(x) \ge 0\) 的情况\begin{enumerate}[label=(\alph*)]
            \item 当 \(x < 0\) 时，\(f(x)=1+x\ge 0 \Rightarrow -1 \leq x < 0\)。
            \item 当 \(x \ge 0\) 时，\(f(x)=1\ge 0 \Rightarrow x \ge 0\)。
        \end{enumerate}$f[g(x)]=1,\ x\ge -1$
    \end{enumerate}
    综上所述，$f[f(x)]=
    \begin{cases}
        2+x&,\ x <-1\\
        1&,\ x \ge -1
    \end{cases}$\par
    \multiplesolution{解二：图解法}
    $f(x)$的图像，如图\ref{fig:ch01-1-2-tx-01-04-01}所示。因此:
        \begin{align*}
        &f(f(x))=\begin{cases}
            1+f(x)&,\ f(x) <0\\
            1&,\ f(x) \ge 0
        \end{cases}\\
        &=\begin{cases}
            2+x&,x <-1\\
            1&,\ x \ge -1
        \end{cases}
    \end{align*}
    \begin{figure}[htbp]
        \centering
        \includegraphics[width=0.8\textwidth]{figures/ch01/1-2/tx-01-04-01.pdf}
        \caption{题型1第【4】题-分段函数 $f(x)$ 图像}
        \label{fig:ch01-1-2-tx-01-04-01}
    \end{figure}
\end{Answer}
\begin{Problem}
    \difficulty{2}
    设
    \(
    f(x)=\begin{cases}
    e^x & x<1\\
    x^2-1 & x\ge 1
    \end{cases},
    \quad
    g(x)=\begin{cases}
    x+2 & x<0\\
    x^2-1 & x\ge 0
    \end{cases},
    \)
    求 $f[g(x)]$.
    \exercisesource{知乎文章：https://zhuanlan.zhihu.com/p/35607988}
    \exercisetip{带入内层函数时需注意；解内层函数不等式需注意}
\end{Problem}
\begin{Answer}
    \multi{解法一：代数法}
    $
    f[g(x)]=\begin{cases}
        e^{g(x)}&,\ g(x) < 1\\
        [g(x)]^{2}-1&,\ g(x) \ge 1
    \end{cases}
    $
    \begin{enumerate}[label=(\arabic*)]
        \item 考虑$
            g(x) <1
        $的情况 \begin{enumerate}[label=(\alph*)]
            \item 当$x<0$时，$g(x)=x+2 <1 \Rightarrow x < -1$
            \item 当$x\ge 0$时，$g(x)=x^2 - 1 <1\Rightarrow 0 \le x <\sqrt{2} $
        \end{enumerate}
        故，$f[g(x)]=e^{g(x)}=\begin{cases}
            e^{x+2}&,\ x < -1\\
            e^{x^2-1} &,0 \le x < \sqrt{2}
        \end{cases}$
        \item 考虑 $
            g(x) \ge 1
        $的情况 \begin{enumerate}[label=(\alph*)]
            \item 当$x < 0$时，$g(x)=x+2 \ge 1 \Rightarrow -1 \le x < 0$
            \item 当$x \ge 0$时，$g(x)=x^2 -1 \ge 1 \Rightarrow x \ge \sqrt{2}$
        \end{enumerate}
        故，$
            f[g(x)]=\begin{cases}
                (x+2)^2-1&,\ -1 \le x <0\\
                (x^2-1)^2-1 &,\ x\ge \sqrt{2}
            \end{cases}
        $
    \end{enumerate} 
    综合(1)(2)所述，得 $
        f[g(x)]=\begin{cases}
            e^{x+2}&,\ x < -1\\
            (x+2)^2-1&,\ -1 \le x < 0\\
            e^{x^2-1}&, \ 0 \le x < \sqrt{2}\\
            (x^2-1)^2-1&,\ x \ge \sqrt{2}
        \end{cases}
    $
    \multi{解法二：图解法}
    略
\end{Answer}
\begin{Problem}
    \difficulty{2}
    设 $
        f(x)=\begin{cases}
            \ln \sqrt{x}&,\ x \ge 1\\
            2x-1&,\ x <1
        \end{cases}
    $求$f[f(x)]$
    \exercisesource{知乎文章：https://zhuanlan.zhihu.com/p/35608680}
    \exercisetip{解内层函数不等式和带入内层函数时容易出错}
\end{Problem}
\begin{Answer}
    \multi{解法一：图解法}
    $f(x)$的图像如图\ref{fig:ch01-1-2-tx-01-05-01}所示。因此，
    \begin{align*}
        &f[f(x)]=\begin{cases}
            \ln\sqrt{f(x)}&,\ f(x) \ge 1\\
            2f(x) -1&,\ f(x) < 1
        \end{cases}\\
        &=\begin{cases}
            \ln\sqrt{\ln \sqrt{x}}&,\ x \ge e^2\\
            4x-3&,\ x < 1\\
            2\ln\sqrt{x}-1&,\ 1 \le x < e^2
        \end{cases}
    \end{align*}
    \begin{figure}[htbp]
        \centering
        \begin{includegraphics}[width =0.8 \textwidth]
            {figures/ch01/1-2/tx-01-05-01.pdf}
        \end{includegraphics}
        \caption{题型1第【5】题-分段函数 $f(x)$ 图像}
        \label{fig:ch01-1-2-tx-01-05-01}
    \end{figure}
\end{Answer}
\begin{Problem}
    \difficulty{4}
    已知 $
        f(x)=\begin{cases}
            \ln x&,\ x \le 0,\\
            x&,\ x >0,
        \end{cases}g(x)=\begin{cases}
            x^2&,\ x \le 1,\\
            x^3&,\ x > 1,
        \end{cases}
    $求$f[g(x)]$
\end{Problem}
\begin{Answer}
    $g(x)$的图像如图\ref{fig:ch01-1-2tx-01-07-01}所示。因此有：
    \begin{figure}[htbp]
        \centering
        \includegraphics[width=0.8 \textwidth]{figures/ch01/1-2/tx-01-07-01.pdf}
        \caption{题型1-【7】题$g(x)$的图像}
        \label{fig:ch01-1-2tx-01-07-01}
    \end{figure}
    \multi{解法一：先内后外}
    \begin{align*}
        f[g(x)]&=\begin{cases}
            f(x^2),\ x \le 1\\
            f(x^3),\ x > 1
        \end{cases}=\begin{cases}
            &\begin{cases}
                ln x^2&,\ x^2>0,\\
                x^2&,\ x^2 \le 0,
            \end{cases}\ x\le 1;\\
            &\begin{cases}
                ln x^3&,\ x^3>0,\\
                x^3&,\ x^3 \le0
            \end{cases}\ x > 1.\\
        \end{cases}\\
        &=\begin{cases}
            \ln x^2&,\ x\le 1\text{并且}x^2>0\\
            x^2&,\ x \le1 \text{并且}x^2\le0\\
            ln x^3&,\ x >1 \text{并且}x^3 >0\\
            x^3&, x >1 \text{并且}x^3\le 0
        \end{cases}=\begin{cases}
            \ln x^2&,\ x\le 1 \text{并}x\neq 0\\
            x^2&,\ x = 0\\
            \ln x^3&,\ x>1\\
        \end{cases}
    \end{align*}
    \multi{解法二：先外后内}
    \begin{align*}
        f[g(x)]&=\begin{cases}
            \ln g(x)&,\ g(x)>0\\
            g(x)&,\ g(x) \le 0
        \end{cases}=\begin{cases}
            \ln x^2&,\ x \le 1\text{且}g(x)>0,\\
            \ln x^3&,\ x > 1\text{且}g(x)>0,\\
            x^3&,\ x\le 1\text{且} g(x)\le 0,\\
            x^2&,\ x>1 \text{且} g(x)\le 0,\\
        \end{cases}\\
        &= \begin{cases}
            \ln x^2&,\ x \neq 0 \text{且}x \le 1\\
            \ln x^3&,\ x >1\\
            x^2&,\ x = 0
        \end{cases}
    \end{align*}
\end{Answer}
\begin{Problem}
    \difficulty{3}
        设 \( g(x)=\begin{cases} 2-x, & x \leq 0, \\ x+2, & x>0, \end{cases} \)
    \( f(x)=\begin{cases} x^2, & x<0, \\ -x, & x \geq 0, \end{cases} \)
    则 \( g[f(x)]=\)(\phantom{5em})
    \begin{enumerate}[label=(\Alph*)]
        \item \( \begin{cases} 2+x^2, & x<0, \\ 2-x, & x \geq 0 \end{cases} \)
        \item \( \begin{cases} 2-x^2, & x<0, \\ 2+x, & x \geq 0 \end{cases} \)
        \item \( \begin{cases} 2-x^2, & x<0, \\ 2-x, & x \geq 0 \end{cases} \)
        \item \( \begin{cases} 2+x^2, & x<0, \\ 2+x, & x \geq 0 \end{cases} \)
    \end{enumerate}
    \exercisesource{真题·1992·数二}
\end{Problem}
\begin{Answer}
    \multi{解法一·先内后外}
    \[
    \begin{aligned}
        &g[f(x)]=\begin{cases}
            g(x^2)&,\ x<0\\
            g(-x)&,\ x \ge 0
        \end{cases}\\
        &=\begin{cases}
            \begin{cases}
                2-x^2,\ x^2 \le 0\\
                x^2+2,\ x^2>0
            \end{cases}&,\ x<0\\
            \begin{cases}
                2-(-x),\ x<0\\
                2+x,\ x \ge 0
            \end{cases}&,\ x \ge 0
        \end{cases}\\
        &=\begin{cases}
            x^2+2,\ x<0\\
            2+x, x \ge 0
        \end{cases}\\
        &\Rightarrow \text{选}D
    \end{aligned}
    \]
    \multi{解法二·先外后内}
    \[
    \begin{aligned}
        &g[f(x)]=\begin{cases}
            2-f(x)&,\ f(x) \le 0\\
            f(x)+2&,\ f(x) > 0
        \end{cases}\\
        &=\begin{cases}
            2-x^2&,\ x<0\text{且}f(x)\le 0\\
            2-(-x)&,\ x \ge 0 \text{且}\le 0\\
            x^2+2&,\ x<0\text{且}f(x)>0\\
            (-x)+2&,\ x \ge 0 \text{且}f(x)>0
        \end{cases}
    \end{aligned}
    \]
\end{Answer}
\begin{Problem}
    \difficulty{2}
    设
    \(
    f(x)=
    \begin{cases}
    1, & |x|\le 1,\\
    0, & |x|> 1,
    \end{cases}
    \)
    则 \(f\{f[f(x)]\}=\)(\phantom{5em})\par
    % 用tabular实现每行2个选项，和图片布局完全一致
    \begin{tabular}{l@{\hspace{20em}}l}
        (A) $0$ & (B) $\begin{cases}0,\ |x| \le 1\\1,\ |x| > 1\end{cases}$ \\[1em]
        (C) $1$ & (D) $\begin{cases}1,\ |x| \le 1\\0,\ |x| > 1\end{cases}$ 
    \end{tabular}
    \exercisesource{b站：BV1rH6kYFEB8}
\end{Problem}
\begin{Answer}
    \multi{解一·先内后外}
    \[
        \begin{aligned}
            f\{f[f(x)]\}&=\begin{cases}
                f\{f[1]\},\ |x| \le 1\\
                f\{f[0]\},\ |x| >1
            \end{cases}\\
            &=\begin{cases}
                f\{1\},\ |x| \le 1\\
                f\{1\},\ |x| >1
            \end{cases}\\
            &=f\{1\}=1
        \end{aligned}
    \]
    \multi{解二·先外后内}
    \[
       \begin{aligned}
            f[f(x)]&=\begin{cases}
                1,\ |f(x)| \le 1\\
                0,\ |f(x)| >1
            \end{cases}\\
            &=\begin{cases}
                1,\ |f(x)| \le 1 \text{并} |x| \le 1\\
                1,\ |f(x)| \le 1 \text{并} |x| >1\\
                0,\ |f(x)| >1 \text{并} |x| \le 1\\
                0,\ |f(x)| >1 \text{并} |x| >1
            \end{cases}\\
            &=\begin{cases}
                1,\ 1 \le 1 \text{并} |x| \le 1\\
                1,\ 0 \le 1 \text{并} |x| >1\\
                0,\ 1 >1 \text{并} |x| \le 1\\
                0,\ 0 >1 \text{并} |x| >1
            \end{cases}\\
            &= \begin{cases}
                1,\ |x| \le 1\\
                1,\ |x| >1
            \end{cases}=1
       \end{aligned}
    \]
    将$f[f(x)]$带入得，\[\begin{aligned}
            f\{f[f(x)]\}&= f\{1\}=1\\
            \Rightarrow \text{选B}
        \end{aligned}
    \]
\end{Answer}
\begin{Problem}
    \difficulty{3}
    设 $
        f(x)= 2x + \sqrt{x^2 + 2x + 1}
    $,\phantom{1em}$g(x)= \begin{cases}
        x + 2,\ x \ge 0\\
        x -1,\ x <0
    \end{cases} $，则$\lim_{x \rightarrow -\frac{1}{3}}g[f(x)]=$(\phantom{5em})
    \exercisesource{b站：BV1rH6kYFEB8·1000题}
    \exercisetip{内层函数不等式的求解}
\end{Problem}
\begin{Answer}
    \multi{解一·先内后外法}
    \[\begin{aligned}
        g[f(x)]&=g[2x+\sqrt{x^2+2x+1}]\\
        &=\begin{cases}
            (2x+\sqrt{x^2+2x+1})+2,\ 2x+\sqrt{x^2+2x+1} \geq 0\\
            (2x+\sqrt{x^2+2x+1})-1,\ 2x+\sqrt{x^2+2x+1} < 0
        \end{cases}\\
        &= \begin{cases}
            2(x+1)+|x+1|,\ x \ge \frac{-1}{3}\\
            2x-1+ |x+1|,\ x < \frac{-1}{3}
        \end{cases}\\
        &= \begin{cases}
            3(x+1),\ x \ge \frac{-1}{3}\\
            x-2,\ x <\frac{-1}{3}
        \end{cases}
    \end{aligned}
    \]
    因此，$\lim_{x\rightarrow -\frac{1}{3}^-}\ g[f(x)] =\lim_{x\rightarrow -\frac{1}{3}^-}\ x-2=-\frac{5}{3}  $\par
    但是$\lim_{x\rightarrow -\frac{1}{3}^+}\ g[f(x)] =\lim_{x\rightarrow -\frac{1}{3}^+}\ 3(x+1)=2 $\par
    所以极限不存在
    \multi{解二·先外后内}
    \[\begin{aligned}
        g[f(x)]&=\begin{cases}
            f(x)+2,\ f(x) \ge 0\\
            f(x)-1,\ f(x) < 0
        \end{cases}\\
        &=\begin{cases}
            2x+2+|x+1|,\ x \ge \frac{-1}{3}\\
            2x-1+|x+1|,\ x < \frac{-1}{3}
        \end{cases}\\
        &=\begin{cases}
            3(x+1),\ x \ge \frac{-1}{3}\\
            x-2,\ x < \frac{-1}{3}
        \end{cases}
    \end{aligned}
    \]
    后续同解法一。
\end{Answer}
\pt{函数奇偶性、周期性、有界性判定}
\pt{分段函数定义域与值域求解}
\pt{函数性质的综合应用}
\pt{应用问题的函数关系建立}
\pt{抽象函数的性质分析}
\pt{函数图像的识别与绘制（结合初等函数性质）}

\pt{复合函数的解析式求解}
\begin{Problem}
    \difficulty{2}
    设 $
        f(x+\dfrac{1}{x})=x^2+\dfrac{1}{x^2}
    $，则$\lim_{x\rightarrow3}\ f(x)=\underline{\quad\quad\quad}$.
    \exercisesource{2026·660·第3题}
    \exercisetip{“$3^2-2$”容易算错}
\end{Problem}
\begin{Answer}
    $
        f(x+\dfrac{1}{x})= (x+\dfrac{1}{x})^2-2 \Rightarrow f(t)=t^2-2
    $\par
    故 $
        \lim_{x\rightarrow3}\ f(x) = \lim_{x\rightarrow3}\ (x^2-2)=3^2-2=7
    $
\end{Answer}
\begin{Problem}
    \difficulty{3}
     已知 \( f\left(\dfrac{1}{x}-1\right)=\dfrac{3x+1}{2x-1} \)，则 \( f(x)=\)(\phantom{5em})
    \begin{enumerate}[label=(\Alph*),leftmargin=2em]
            \item \( \dfrac{4-x}{1+x} \)
            \item \( \dfrac{4+x}{1-x} \)
            \item \( \dfrac{1-x}{4+x} \)
            \item \( \dfrac{3x+1}{2x-1} \)
    \end{enumerate}
    \exercisesource{b站：BV1rH6kYFEB8-给出$f(\text{狗})$求$f(x)$}
\end{Problem}
\begin{Answer}
    令 $
        t=\dfrac{1}{x}-1
    $，则$\dfrac{1}{x} = 1+t$\par
    则 $
        f(t)=\dfrac{3+\dfrac{1}{x}}{2-\dfrac{1}{x}}=\dfrac{3+(1+t)}{2-(1+t)}=\dfrac{4+t}{1-t}
    $\par
    所以，$
        f(x)=\dfrac{4+x}{1-x}\Rightarrow
    $答案选$B$
\end{Answer}
\begin{Problem}
    \difficulty{3}
    已知$f\left(x+\dfrac{1}{x}\right)=\dfrac{x+x^3}{1+x^4}$，求$f(x)$
    \exercisesource{b站：BV1rH6kYFEB8-给出$f(\text{狗})$求$f(x)$}
\end{Problem}
\begin{Answer}
    令 $
        t = x + \dfrac{1}{x}
    $，则 $
        t^2 = x^2+\dfrac{1}{x^2}+2
    $\par
    $
        f(t) \upeq{分子分母同时除以$x^2$}{6em} \dfrac{\dfrac{1}{x}+x}{\dfrac{1}{x2}+x^2} = \dfrac{t}{t^2-2} \Rightarrow f(x)=\dfrac{x}{x^2-2}
    $
\end{Answer}
\begin{Problem}
    \difficulty{3}

    \exercisesource{b站：BV1rH6kYFEB8-给出$f(\text{狗})$求$f(x)$}
\end{Problem}
 \begin{Answer}
    
 \end{Answer}
